% Section 1: Introduction.
\section{Introduction}\label{sec:intro}

\subsection{Setting}\label{sec:setting}
Let $p>0$, $s>1$, and let $\Lambda=\{\lambda_m\}_{m\ge1}$ be an increasing sequence of positive numbers with
$\lambda_m\sim m^p$. We call $\Lambda$ a \emph{node set} and study the discrete measure on its reciprocals
\begin{equation}\label{eq:sigmaL}
  \sigma_\Lambda:=\sum_{m\ge1}\lambda_m^{-s/p}\,\delta_{1/\lambda_m}.
\end{equation}
For the \emph{lattice} $\lambda_m=m^p$ this is the measure $\sigma_{p,s}=\sum_kk^{-s}\delta_{k^{-p}}$ of the earlier
papers of this series \cite{PaperI,PaperII,Paper3}. Let $P_n$ be the monic orthogonal polynomials of $\sigma_\Lambda$;
they satisfy
\begin{equation}\label{eq:rec}
  xP_n=P_{n+1}+a_nP_n+b_n^2P_{n-1}\qquad(n\ge0,\ b_0:=0,\ b_n>0),
\end{equation}
and $a_n,b_n$ are the entries of the Jacobi matrix whose spectrum is $\{1/\lambda_m\}\cup\{0\}$. With
\begin{equation}\label{eq:cp}
  I(p):=\frac{\sqrt\pi}{2}\,\frac{\Gamma(\frac1p+\frac12)}{\Gamma(\frac1p+1)},\qquad
  c(p):=\frac14\Bigl(\frac{\pi}{2I(p)}\Bigr)^{p},
\end{equation}
the lattice satisfies the pointwise law $n^pb_n\to c(p)$, $n^pa_n\to2c(p)$, with rate $O(1/n)$ and an explicit $1/n$
term for every $p>0$ \cite[Theorems~1.2 and~1.4]{PaperII}. We write $c=c(p)$,
\begin{equation}\label{eq:alpha-y}
  \alpha:=\frac{s-1-p}{p},\qquad y_n:=\frac{n^pb_n}{c}-1,
\end{equation}
so that $\alpha>-1$ is the Bessel order of the hard edge at $0$ in \cite{PaperII}. The questions of this paper are:
how does the expansion of $y_n$ continue, and how does it depend on $\Lambda$? The organising object is the
\emph{node zeta function}
\begin{equation}\label{eq:Z}
  Z_\Lambda(w):=\sum_{m\ge1}\lambda_m^{-w},
\end{equation}
continued meromorphically from $\Re w>1/p$; for the lattice $Z_\Lambda(w)=\zeta(pw)$. Two quantities recur: the
numbers
\begin{equation}\label{eq:W}
  W(q):=\binom{2q}{q}c^q=\frac{(4c)^q\,\Gamma(q+\frac12)}{\sqrt\pi\,\Gamma(q+1)}\qquad(q\in\C\setminus\{-\tfrac12,-\tfrac32,\dots\}),
\end{equation}
which satisfy $W(1/p)=1$ for every $p$ (this is the definition \eqref{eq:cp} of $c$), and the coefficients
\begin{equation}\label{eq:epsj}
  \varepsilon_j:=W(j)\,p\,\zeta(-jp)=\binom{2j}{j}p\,c^j\zeta(-jp)\qquad(j\ge1),
\end{equation}
so that $\varepsilon_1=\kp:=2pc\,\zeta(-p)$ is the constant of \cite[Theorem~1.4(ii)]{PaperII}.

\subsection{Results}\label{sec:results}
The first result is an expansion of the lattice coefficients far beyond the $1/n$ term. Its form is simplest in the
shifted index
\begin{equation}\label{eq:N}
  N:=n+\frac\alpha2+\frac12 .
\end{equation}
The rational functions $f_4,f_6$ of $(p,\alpha)$ and $g_{1,4}$ are given in \eqref{eq:f4}--\eqref{eq:g14}.

\begin{mainthm}[the lattice to order $n^{-7}+n^{-2-3p}$]\label{thm:A}
Let $p>0$, $p\ne1$, $s>1$. The recurrence coefficients of $\sigma_{p,s}$ satisfy
\begin{multline}\label{eq:thmA}
  \frac{n^pb_n}{c}=\Bigl(\frac nN\Bigr)^p\Bigl[1+\frac{f_2}{N^2}+\frac{f_4}{N^4}+\frac{f_6}{N^6}
  +\frac{\varepsilon_1}{N^{1+p}}\Bigl(1+\frac{g_1}{N^2}+\frac{g_{1,4}}{N^4}\Bigr)\\
  +\sum_{j\ge2}\frac{\varepsilon_j}{N^{1+jp}}\Bigl(1+\frac{g_j}{N^2}\Bigr)
  +\frac{3p+1}{2p}\,\frac{\varepsilon_1^2}{N^{2+2p}}\Bigr]
  +O\bigl(n^{-7}+n^{-2-3p}\bigr),
\end{multline}
where the sum is over the finitely many $j\ge2$ with $1+jp<\min(7,2+3p)$,
\[
  f_2=\frac{p+1}{12}-\frac{(p+2)\alpha^2}{8},\qquad g_j=A_j+B_j\alpha^2,\qquad
  B_j=\frac{(p+2)\bigl(j^2p^2+(j+1)p+2\bigr)}{8(2j-1)p},
\]
\[
  A_j=\frac{j^2(j^2+2j-2)p^3+2(j+1)(2j^2-1)p^2+(5j^2+2j-6)p+2(j-2)}{24(2j-1)p}.
\]
\end{mainthm}

Theorem~\ref{thm:A} is the case $K=6$, $B=2$ of Theorem~\ref{thm:N6}, which gives a double expansion of
$n^pb_n/c$ in $1/n$ and in the parameters $\mu_{n,j}=2\zeta(-jp)/(jn^{jp})$ of the band factor, to every order, through
a finite residue recursion. Its proof combines the third-order Riemann--Hilbert analysis of \cite{PaperII}, extended
to order $K$, with exact computer algebra over $\mathbb Q(p,\alpha)$ and, for the zeta terms, over
$\mathbb Q(p,\alpha,j)$ with $j$ a symbol. Three structural facts are visible in~\eqref{eq:thmA}: the shift
$n\mapsto N$ absorbs the whole $1/n$ term and the $(s-1)^2$ part of the $n^{-2}$ term; every layer computed is even in
$1/N$; and the first-order zeta terms are dressed by an even factor $1+g_jN^{-2}+\dotsb$ with $g_j$ quadratic in $\alpha$.
Not every displayed term is a coefficient of $b_n$ for every $p$: Table~\ref{tab:visible} lists which are.

The second result explains the zeta coefficients. The band factor of \cite{PaperII} enters the first-order
correction through one linear functional $e_n$ (Section~\ref{sec:kernel}), whose asymptotic series is
$\sum_j\varepsilon_jn^{-jp}$.

\begin{mainthm}[the zeta layer as a Mellin--Barnes integral]\label{thm:B}
Let $p>0$, $p\ne1$, and $0<c_0<1$. The band-factor functional $e_n$ of the lattice satisfies
\[
  e_n=\frac1{2\pi i}\int_{c_0-i\infty}^{c_0+i\infty}p\,W(w)\,\zeta(-wp)\,n^{-wp}\,
  \pi\Bigl[\cot\frac{\pi wp}2-\cot(\pi w)\Bigr]\diff w ,
\]
and, for every $J\ge1$, $e_n=\sum_{j\le J}\varepsilon_jn^{-jp}+O(n^{-(J+1)p})$. The kernel $-\pi\cot(\pi w)W(w)(4c)^{-w}$
is the Mellin transform of the truncated per-node law $(1-u)^{-1/2}\mathbf 1_{u<1}$ (Lemma~\ref{lem:K0}).
\end{mainthm}

The third result is on the side of potential theory. In the variable $x=n^p\cdot(\text{node})$, the $n$ nodes that
carry the polynomial $P_n$ have the continuum density $\sigma(y)=\frac1py^{-1/p-1}$, and the zero distribution of
$P_n$ is the equilibrium measure of unit mass constrained by $\sigma$ \cite{Rakhmanov,DragnevSaff,PaperI}. We solve
this problem for every constraint with $y\sigma(y)$ decreasing.

\begin{mainthm}[constrained equilibrium]\label{thm:C}
Let $\sigma$ be a positive, locally H\"older continuous density on $(0,Y]$, $0<Y\le\infty$, with $y\sigma(y)$ strictly
decreasing and, if $Y=\infty$, $\int_1^\infty\sigma(y)\log y\diff y<\infty$. Then the soft-edge condition
$\int_a^Y\sigma(y)(1-a/y)^{-1/2}\diff y=1$ has exactly one root $a\in(0,Y)$, and the constrained equilibrium measure is
$\nu\diff x|_{[0,a]}+\sigma\diff y|_{[a,Y]}$ with the explicit density $\nu$ of~\eqref{eq:nu}, $0<\nu<\sigma$. For the
lattice $a=4c(p)$. For a counting function $\mathcal N(\lambda)\sim\sum_\rho c_\rho\lambda^\rho$, the root
$a_n=4c(1+\eta_n)n^{-p}$ of $\int(1-a_n\lambda)^{-1/2}\mathbf 1_{a_n\lambda<1}\diff\mathcal N(\lambda)=n$ satisfies
\[
  n\eta_n=p\sum_{j\ge0}W(j)Z_{\mathcal N}(-j)n^{-jp}+p\sum_qR_qW(q)\,\pi\cot(\pi q)\,n^{-pq}+O(n^2\eta_n^2),
\]
the second sum over the poles $q\ne-1/p$ of $Z_{\mathcal N}(-w)=\int\lambda^{w}\diff\mathcal N(\lambda)$, with residues
$R_q$ (Proposition~\ref{prop:M12}).
\end{mainthm}

The equilibrium statement holds up to one input from potential theory, the positivity of the logarithmic energy of
compactly supported signed measures of mass zero (Lemma~\ref{lem:energy}), which we use without proof.

\subsection{The spectral dictionary}\label{sec:dictionary}
For the lattice, Theorems~\ref{thm:A} and~\ref{thm:B} identify the coefficients of $b_n$ with values of $\zeta$. For a
general node set the same Riemann--Hilbert analysis has not been carried out, and we formulate what the theorems and
the numerical evidence suggest (Section~\ref{sec:conj}). In the natural index $N_\Lambda:=n+\frac\alpha2-Z_\Lambda(0)$
(which is~\eqref{eq:N} for the lattice, $\zeta(0)=-\frac12$), the expansion of $n^pb_n/c$ should read, at leading
order in each spectral datum,
\[
  \Bigl(\frac n{N_\Lambda}\Bigr)^{\!p}\Bigl[1+\frac{f_2}{N_\Lambda^2}+\dotsb
  +\sum_{j\ge1}\frac{W(j)\,p\,Z_\Lambda(-j)}{N_\Lambda^{1+jp}}
  +\sum_q\frac{p\,R_q\,W(q)\,\pi\cot(\pi q)}{N_\Lambda^{1+pq}}+\dotsb\Bigr],
\]
the second sum running over the poles $q$ of $Z_\Lambda(-w)-\zeta(-pw)$, with $w$-residues $R_q$. The integer residues
of $\pi\cot(\pi w)W(w)pZ_\Lambda(-w)$ are the zeta terms of Theorem~\ref{thm:B}; its residues at the poles of $Z_\Lambda$
are the pole terms of Theorem~\ref{thm:C}. The dictionary is exact for the measures on reciprocal squared Bessel zeros
(Proposition~\ref{prop:bessel}), it reads the index offset of the spectrum from $Z_\Lambda(0)$, and it
predicts, for the eigenvalues of $-\frac{\diff^2}{\diff q^2}+q^N$ with $N$ even, an $n^{-2}$ shift equal to
$(N-1)/(6(N+2))$ (Section~\ref{sec:quartic}). For drums at $p=1$, outside the range of the theorems, the perimeter
appears at order $n^{-1/2}$ and the corners at order $n^{-1}$, as the kernel predicts (Section~\ref{sec:drums}). The
dictionary has a limit: subleading terms of each pole, one order lower, are not given by it
(Remark~\ref{rem:sublead}), and for spectra whose nodes deviate from a smooth node function by a fraction of the
spacing the expansion is replaced by noise of relative size $n^{-1/2}$, whose variance appears numerically to be set by the
spectral form factor (Section~\ref{sec:noise}).

\subsection{Organisation and notation}\label{sec:organisation}
Section~\ref{sec:prior} restates the results of the earlier papers that are used, together with two results of the
same project that extend them by one order. Section~\ref{sec:equilibrium} proves Theorem~\ref{thm:C},
Section~\ref{sec:kernel} proves Theorem~\ref{thm:B}, and Section~\ref{sec:natural} proves Theorem~\ref{thm:A} and its
general form. Section~\ref{sec:examples} collects the examples and the numerical evidence, and
Section~\ref{sec:conj} the conjectures. Section~\ref{sec:literature} will relate the results to the literature.
Appendix~\ref{app:repro} lists the scripts.

Throughout, $p>0$ and $c=c(p)$; $\alpha$, $y_n$, $W$, $\varepsilon_j$, $N$ are as in \eqref{eq:alpha-y}--\eqref{eq:N}.
The Bessel order $\alpha$ is that of \cite{PaperII}, and $\kp=\varepsilon_1$; the constant $\kappa_p$ of
\cite{PaperI} does not occur. We write $\binom{2q}{q}=\Gamma(2q+1)/\Gamma(q+1)^2$ for complex $q$. The letter $N$
denotes the natural index, except in Section~\ref{sec:quartic}, where $q^N$ is a potential. Constants $C$ may change
from line to line. Numbers obtained by fitting floating-point data are quoted with the spread of the fit over
fit windows and nuisance exponents in parentheses; such a spread is not an error bound.
